\chapter{Antecedentes}\label{cap3}

En la Unidad Interdisciplinaria de Ingeniería Campus Zacatecas (UPIIZ-IPN), se
han realizado diversas investigaciones relacionadas con la presente propuesta. A
continuación, se muestran los trabajos más relevantes que contribuyen al desarrollo
de los objetivos planteados.

En 2016, los alumnos Pedro Rodríguez Segura y Juan José Saldívar Nava,
asesorados por el M.\ en C.\ Omar Désiga Orenday, Dra.\ Verónica Ávila Vázquez y
el M.\ en C.\ Sergio Zavala Castillo, presentaron el proyecto ``Reactor para la
producción de biodiesel utilizando aceite vegetal usado''. Este trabajo se enfocó en
el control de variables para la producción de biodiesel como la temperatura,
velocidad de agitación e intensidad luminosa \cite{Rodriguez2016}. Su utilidad radica en que
proporciona bases importantes para la implementación de sistemas de medición y
control de temperatura y agitación en medios líquidos.

Por otro lado, en 2023, el trabajo ``Diseño y construcción de un fotobiorreactor como
soporte vital para prueba en la estratosfera'' desarrollado por los alumnos Carlos
Alvarado, Joel Escareño y Diego Osorio, asesorados por el Dr.\ Hans Christian
Correa Aguado, M.\ en C.\ Flabio Darío Mirelez Delgado y M.\ en I.\ Addried Samir
Moreno Castro, tuvo como fin controlar las variables que afectan al desarrollo del
crecimiento de las microalgas, siendo estas la intensidad lumínica y la temperatura,
entre otros \cite{Alvarado2023}. Este proyecto será de gran importancia ya que sienta las bases para
el diseño de sistemas de monitoreo y control de variables.
